\documentclass[11pt]{article}

% ---------------------------------------------------------------- packages
\usepackage[margin=1in]{geometry}
\usepackage{amsmath,amssymb,amsthm}
\usepackage{bm}
\usepackage{booktabs}
\usepackage{multirow}
\usepackage{graphicx}
\usepackage[authoryear,sort&compress]{natbib}
\usepackage[colorlinks=true,linkcolor=blue,citecolor=blue,urlcolor=blue]{hyperref}
\usepackage{microtype}
\emergencystretch=1.5em

% ---------------------------------------------------------------- macros
\newcommand{\R}{\mathbb{R}}
\newcommand{\D}{\mathcal{D}}
\newcommand{\V}{\mathcal{V}}
\newcommand{\Vsh}{\mathcal{V}_{\mathrm{sh}}}
\newcommand{\E}{\mathbb{E}}
\newcommand{\norm}[1]{\left\lVert #1 \right\rVert}
\newcommand{\bhat}[1]{\hat{\bm{#1}}}
\newcommand{\argmin}{\operatorname*{arg\,min}}
\newcommand{\argmax}{\operatorname*{arg\,max}}
\newcommand{\softmax}{\operatorname{softmax}}
\newcommand{\KL}{\operatorname{KL}}
\newcommand{\todo}[1]{\textcolor{red}{[TODO: #1]}}
\theoremstyle{definition}
\newtheorem{proposition}{Proposition}
\newtheorem{remark}{Remark}

\title{\bf Fingerprinting Language-Model Weights via the Distributional Divergence of\\ Aligned Representation-Equivalent Probes}

\author{Hanlin Tan}

\date{}

\begin{document}
\maketitle

\begin{abstract}
Determining whether two language models have (nearly) identical weights---a
question at the heart of model-copy detection, provenance attribution, and
adversarial weight-tampering auditing---is made difficult by the inaccessibility
of hidden states in production settings and by the fact that two models may
produce identical outputs while differing in their parameters. We introduce a
two-stage approach that attacks this problem from a purely observable,
distribution-level signal. First, we generate a set of
\emph{representation-equivalent probes}: sentences that the model treats as
nearly indistinguishable from a reference sentence---their pooled intermediate
representation is forced into a tight cluster around the reference's---yet
differ from it on the surface. Equivalence here is defined relative to the
model's own representation, not to human semantics: the probes need not be
synonyms in any linguistic sense. They are produced by a HotFlip-style discrete
optimization restricted to a shared subword vocabulary. Second, we treat the induced next-token (log)
probability distributions of these probes as the observable signal and measure
their per-probe \emph{distributional divergence} (KL) as a fingerprint of the
model. We show that this fingerprint is a monotonic, sensitive detector of
weight perturbation: adding zero-mean Gaussian noise of scale $r\sigma$ to the
parameters of three open-weight models (Qwen3.5-0.8B, Llama-3.2-1B,
Gemma-3-1B) yields calibration curves whose divergence grows quadratically in
$r$ (as the Fisher-information predicts), stable across model families,
intermediate layers, and vocabulary restrictions. We further demonstrate a
zero-noise change detector over the
DeepSeek API's top-$20$ log-probability surface, and identify the tokenizer
mismatch between distinct model families as the principal obstacle to
cross-family comparison. Our results establish distributional divergence as a
practical, black-box-compatible fingerprint for weight equivalence, while
delineating the conditions under which it is---and is not---sufficient.
\end{abstract}

\section{Introduction}
\label{sec:intro}

Large language models are expensive to train, and their weights are the primary
carrier of that investment. This raises a natural and increasingly practical
question: given two models, can we determine whether their parameters are
(nearly) identical? The question matters for at least three reasons. First,
\emph{copy detection}: verifying whether a released checkpoint has been
surreptitiously derived from a proprietary model. Second,
\emph{provenance}: attributing a model to its training lineage. Third,
\emph{tamper auditing}: detecting whether a served model has been silently
modified---for instance, by a malicious actor injecting fine-tuned backdoors or
by an operator swapping weights at a load balancer.

Answering this question is hard for two coupled reasons. On one hand, the
weights themselves are typically not exposed: a deployment only offers token
distributions (or, at most, top-$k$ log-probabilities) through an API. On the
other, \emph{output equivalence does not imply weight equivalence}---two very
different parameterizations can realize the same input--output mapping, so
agreement on a held-out set of prompts is a necessary but weak signal.

We propose to sidestep both difficulties by exploiting the
\emph{second-order} response of the model's observable distribution: the
distributional divergence of the observable next-token distribution, whose
quadratic response is controlled by the Fisher information. The key observation
is that a model's weights induce a specific local geometry in representation
space. If we can produce a set of sentences that are \emph{indistinguishable to
the model}---in the sense that they map to a tight cluster under its own
intermediate representation---then the way in which these
representation-equivalent sentences spread in next-token probability space is a
signature of the weights themselves. Perturbing the weights perturbs this
distributional geometry in a way that a single-output comparison cannot reveal.
Throughout, ``equivalence'' is defined by the model, not by a human reader: the
probes must land on (nearly) the same latent representation, and are
\emph{not} required to be synonyms in any linguistic sense.

Concretely, we make three contributions:

\begin{enumerate}
  \item \textbf{Constrained representation-equivalent generation}
  (\S\ref{sec:syn}): a discrete optimization procedure that, given a
  reference sentence, emits multiple sentences the model treats as nearly
  indistinguishable from it---their intermediate-layer representations are
  aligned to the reference's---while differing on the surface, using only tokens
  from a restricted (shared) subword vocabulary. The output is always a discrete
  token sequence, never a vector.

  \item \textbf{Distributional-divergence fingerprinting}
  (\S\ref{sec:cov}): a fingerprint defined as the mean per-probe next-token KL
  divergence between two models, together with a two-sample change-detection
  test (per-token $\chi^2$ as the primary verdict, plus mean-vector and
  covariance-spectrum diagnostics).

  \item \textbf{Empirical calibration and limits} (\S\ref{sec:exp}): a
  calibration curve showing that the fingerprint diverges monotonically with
  added Gaussian weight noise, stable across layers and vocabulary
  restrictions; a zero-noise change detector over the DeepSeek API; and a
  characterization of the tokenizer-mismatch barrier to cross-family
  comparison.
\end{enumerate}

\section{Related Work}
\label{sec:related}

\paragraph{Discrete adversarial attacks.}
Our representation-equivalent generator descends from the family of gradient-based
discrete text optimizers initiated by \citet{hotflip}, which approximates the
effect of a token substitution by the inner product of the input-embedding
gradient and the candidate embedding. Subsequent work refined this idea with
beam search and constrained candidate sets \citep{textfooler}. We differ in
purpose---we optimize for \emph{preservation of the model's representation},
requiring the output to remain nearly indistinguishable to the model itself
rather than for label flipping---and in the constraint that candidates be drawn
from a shared subword vocabulary, which is essential for cross-model
comparison.

\paragraph{Representation similarity.}
Measuring the similarity of neural representations has a long history, with
centered kernel alignment (CKA) \citep{cka} and representational similarity
analysis \citep{rsa} as the canonical tools. These methods compare the
\emph{internal activations} of two networks and thus require white-box access.
Our fingerprint is instead computed on the observable log-probability surface,
making it compatible with API-only deployments, while still being grounded in
the intermediate representation through the probe-construction step.

\paragraph{Model fingerprinting and extraction.}
Model extraction attacks \citep{modelsteal} and watermarking
\citep{watermark} both fingerprint a model's \emph{behavior}; weight-recovery
attacks \citep{rolnick} recover the \emph{parameters} directly. Our work
occupies the middle ground: we estimate a second-order statistic that is
sensitive to the weights but does not require recovering them, and is
observable without hidden states.

\paragraph{Query-based model fingerprinting.}
The closest prior art infers a model's origin by actively probing it with
optimized queries \citep{teacherfingerprint}. We extend this line of work in
two directions: our probes are \emph{actively synthesized} to be
representation-equivalent---indistinguishable in the model's own intermediate
representation, not in any linguistic sense---so the signal is a
within-representation-cluster distributional divergence rather than an
across-meaning mean; and our statistic is the \emph{KL divergence} between the
per-probe next-token distributions rather than their mean, which turns out to
be far more sensitive to weight perturbation.

\section{Problem Formulation}
\label{sec:form}

Let a language model $M_\theta$ be parameterized by weights $\theta \in \R^P$,
with vocabulary $\V$ and input-embedding matrix $\bm{E} \in \R^{|\V| \times D}$.
For an input token sequence $\bm{x} = (x_1, \dots, x_T) \in \V^T$, denote by
$\bm{h}^{(\ell)}(\bm{x}) \in \R^{T \times D}$ the hidden states at intermediate
layer $\ell$, and by $\bar{\bm{h}}^{(\ell)}(\bm{x}) = \frac1T
\sum_t \bm{h}^{(\ell)}_t(\bm{x}) \in \R^D$ its mean-pooled sentence vector.
The model maps a prefix to a next-token distribution
$p_\theta(\cdot \mid \bm{x}) = \softmax(\bm{W} \bm{h}^{(L)}(\bm{x}))$ over the
vocabulary.

\paragraph{Goal.}
Given black-box access to two models $M_\theta$ and $M_{\theta'}$ (top-$k$
log-probabilities only), decide whether $\norm{\theta - \theta'} \le \epsilon$
for a small $\epsilon$, i.e., whether their weights are nearly identical.

\paragraph{Two-part strategy.}
We first construct a \emph{probe set} $\mathcal{P} = \{\bm{x}^{(1)},
\dots, \bm{x}^{(m)}\}$ of representation-equivalent sentences---sentences the
model's intermediate representation cannot distinguish from a common reference
(\S\ref{sec:syn}). Each probe
$\bm{x}^{(i)}$ induces an observable distribution $p_\theta(\cdot \mid
\bm{x}^{(i)})$; we collect a fingerprint statistic from these distributions
(\S\ref{sec:cov}). Equivalence is then tested by comparing two fingerprints
collected at different times or from different models (\S\ref{sec:detect}).

\section{Method}
\label{sec:method}

The fingerprinting procedure below assumes the model is served through an API
that exposes token-level log-probabilities---specifically, the top-$k$ next-token
log-probabilities at each decoding position. Every subsequent step, from the
representation-equivalent probes of \S\ref{sec:syn} to the distributional-divergence
fingerprint of \S\ref{sec:cov}, reads these log-probabilities directly; a provider that
withholds or quantizes them makes the fingerprint unobservable, as noted in
\S\ref{sec:disc}.

\subsection{Constrained Representation-Equivalent Generation}
\label{sec:syn}

We want, for a reference sentence $\bm{s}$ with latent representation
$\bm{h}^\star = \bar{\bm{h}}^{(\ell)}(\bm{s})$, a collection of sentences
$\{\bm{s}^{(j)}\}$ such that (i) each $\bm{s}^{(j)}$ is
\emph{representation-equivalent} to $\bm{s}$---indistinguishable from it in the
model's own intermediate representation---and (ii) $\bm{s}^{(j)} \ne \bm{s}$ on
the surface, and (iii) every token of $\bm{s}^{(j)}$ lies in a
\emph{restricted subword vocabulary} $\Vsh \subseteq \V$.

\paragraph{Equivalence as latent alignment.}
The notion of equivalence in (i) is deliberately \emph{model-relative}: two
sentences count as equivalent for our purposes exactly when the model's own
intermediate representation places them in a tight cluster. We do not require,
and in general do not obtain, synonymy in the linguistic sense---the pair
\emph{``the lazy dog''} $\to$ \emph{``the lazy cat''} of \S\ref{sec:exp:syn}
is a model-equivalence, not a semantic one. We operationalize (i) as the
requirement that the pooled intermediate representation be close to the
reference:
\begin{equation}
  \bm{s}^{(j)} = \argmin_{\bm{x} \in (\Vsh)^T}
  \ \norm{\bar{\bm{h}}^{(\ell)}(\bm{x}) - \bm{h}^\star}^2
  \ \ \text{s.t.}\ \ \bm{x} \ne \bm{s},
  \label{eq:align}
\end{equation}
Because $\bm{x}$ ranges over a discrete set, we optimize
Eq.~\eqref{eq:align} by first-order gradient guidance.

\paragraph{HotFlip-style search.}
Let $\bm{e}_i \in \R^D$ be the input embedding of token $x_i$, assembled into
$\bm{E}_{\bm{x}} \in \R^{T \times D}$. The loss is
$\mathcal{L}(\bm{x}) = \norm{\bar{\bm{h}}^{(\ell)}(\bm{x}) - \bm{h}^\star}^2$.
A first-order Taylor expansion gives the approximate change in loss when token
$x_i$ is replaced by token $t$ as
\begin{equation}
  \Delta \mathcal{L}(i, t) \ \approx\ \bm{g}_i^\top \bm{e}_t,
  \label{eq:hotflip}
\end{equation}
where $\bm{g}_i = \nabla_{\bm{e}_i} \mathcal{L}$ and $\bm{e}_t$ is the
embedding of candidate token $t$. Ranking all of $\Vsh$ by this score is
expensive, and---because the gradient vanishes at the reference---admits
degenerate, non-fluent substitutions; we therefore \emph{instantiate} the search
with a restricted candidate set. For each position the candidates are the $q$
nearest \emph{embedding neighbors} of the current token within $\Vsh$; neighbors
are surface-different from the reference by construction, so no explicit
diversity term is needed. We then \emph{exactly} re-score every candidate
substitution with a forward pass and accept the one that most reduces the
alignment loss (Algorithm~\ref{alg:syn}).

\begin{figure}[t]
\centering
\fbox{\begin{minipage}{0.94\linewidth}
\smallskip
{\small\ttfamily
Algorithm~1 \ --- \ Embedding-neighbor representation-equivalent generation.\par\smallskip
\hangindent=1.5em \textbf{Input:} reference $\bm{s}$, layer $\ell$, shared
vocab $\Vsh$, neighbor count $q$.\par
\hangindent=1.5em $\bm{h}^\star \gets \bar{\bm{h}}^{(\ell)}(\bm{s})$\par
\hangindent=1.5em $\mathcal{C} \gets \bigcup_i \big\{ (i,t) : t \in \Vsh,\
t \ne s_i,\ t \in \mathrm{NN}_q(s_i) \big\}$\par
\hangindent=1.5em $\bm{x}^\star \gets \argmin_{(i,t) \in \mathcal{C}}\
\norm{\bar{\bm{h}}^{(\ell)}(\bm{x}_{i\to t}) - \bm{h}^\star}^2$\par
\hangindent=1.5em \quad ($\bm{x}_{i\to t}$ denotes $\bm{s}$ with token $i$
replaced by $t$)\par
\hangindent=1.5em \textbf{return} $\bm{x}^\star$ if it improves over $\bm{s}$,
else $\bm{s}$\par
}
\smallskip
\end{minipage}}
\caption{Embedding-neighbor representation-equivalent generation.}
\label{alg:syn}
\end{figure}


\subsection{Distributional-Divergence Fingerprint}
\label{sec:cov}

Given a probe set $\mathcal{P} = \{\bm{x}^{(1)}, \dots, \bm{x}^{(m)}\}$, each
probe induces a next-token distribution over the vocabulary, which we represent
by its log-probability vector restricted to the shared vocabulary $\Vsh$:
\begin{equation}
  \bm{\phi}_\theta(\bm{x}) = \big[\ \log p_\theta(t \mid \bm{x}) \big]_{t \in \Vsh}
  \in \R^{|\Vsh|},
  \label{eq:phi}
\end{equation}
so that $p_\theta(\cdot \mid \bm{x}) = \exp \bm{\phi}_\theta(\bm{x})$ is a
probability distribution over $\Vsh$. In the white-box (local) setting,
$\bm{\phi}$ is available for the full vocabulary; in the black-box (API)
setting, only the top-$k$ entries are observed and the remainder are set to a
fixed floor (a log-probability truncated to $-30$), which approximates the
full-vocabulary distribution well because the probability mass concentrates in
the top-$k$ tokens. To reduce sampling variance we average $\bm{\phi}$ over
$S$ independent draws per probe, yielding per-probe means
$\bar{\bm{\phi}}^{(i)}$.

We define the \emph{fingerprint} of $M_\theta$ as the collection of per-probe
next-token distributions $\{\bar{\bm{\phi}}^{(1)}, \dots,
\bar{\bm{\phi}}^{(m)}\}$, and measure the difference between two models
$\theta$ and $\theta'$ by the mean per-probe KL divergence
\begin{equation}
  D_{\mathrm{KL}}(\theta, \theta')
  \;=\; \frac1m \sum_{i=1}^m
    \KL\!\big(p_\theta(\cdot \mid \bm{x}^{(i)})\ \|\ p_{\theta'}(\cdot \mid \bm{x}^{(i)})\big),
  \label{eq:dkl}
\end{equation}
computed with the observed top-$k$ entries and the floor for the remainder.
The KL divergence is the central object: it measures how far the next-token
distributions induced by sentences that are \emph{representation-equivalent}---
indistinguishable to the model itself---have moved under a change of
parameters, a quantity that is a direct functional of the weights but is not
exposed by any single forward pass.

\subsection{Theoretical Analysis}
\label{sec:theory}

The calibration results of \S\ref{sec:exp:calib} raise two questions that the
definition above does not answer: \emph{why} the fingerprint responds
monotonically---and $D_{\mathrm{KL}}$ approximately quadratically---to Gaussian
weight noise, and \emph{how} this divergence relates to the $\ell_2$ distance
between the weights. Both are made precise by the Jacobian of the fingerprint
with respect to the weights. For a fixed probe $\bm{x}$, define
\begin{equation}
  \bm{J}_\theta(\bm{x}) = \frac{\partial \bm{\phi}_\theta(\bm{x})}{\partial \theta}
    \in \R^{|\Vsh| \times P},
  \label{eq:jacobian}
\end{equation}
which is well defined because $\bm{\phi}_\theta$ is a smooth function of the
logits and the logits are smooth functions of $\theta$. Write
$\bm{\phi}_i = \bm{\phi}_\theta(\bm{x}^{(i)})$, $\bm{J}_i =
\bm{J}_\theta(\bm{x}^{(i)})$, and $\bar{\bm{J}} = \tfrac1m \sum_i \bm{J}_i$.

\begin{proposition}[The Jacobian is a score function]
\label{prop:score}
Let $z_t(\bm{x}) = [\bm{W}\bm{h}^{(L)}(\bm{x})]_t$ be the logit of token $t$.
Then
\begin{equation}
  \big[\bm{J}_\theta(\bm{x})\big]_{t,:} \;=\;
  \nabla_\theta z_t(\bm{x}) - \sum_{t' \in \V} p_\theta(t' \mid \bm{x})\,\nabla_\theta z_{t'}(\bm{x})
  \;=\; \nabla_\theta \log p_\theta(t \mid \bm{x}).
  \label{eq:score}
\end{equation}
\end{proposition}

\begin{proposition}[Quadratic response and the Fisher information]
\label{prop:fisher}
For a fixed probe $\bm{x}$, the KL divergence between the next-token
distributions of $\theta$ and $\theta + \delta$ admits the second-order
expansion
\begin{equation}
  \KL\!\big(p_\theta(\cdot \mid \bm{x})\ \|\ p_{\theta+\delta}(\cdot \mid \bm{x})\big)
  \;=\; \tfrac12\, \delta^\top \bm{F}_\theta(\bm{x})\, \delta
      + O(\norm{\delta}_2^3),
  \label{eq:fisherexpand}
\end{equation}
where the first-order term vanishes identically (since
$\E_{t \sim p_\theta(\cdot \mid \bm{x})}[\nabla_\theta \log p_\theta(t \mid \bm{x})] = 0$)
and
\begin{equation}
  \bm{F}_\theta(\bm{x})
  \;=\; \E_{t \sim p_\theta(\cdot \mid \bm{x})}
      \big[\nabla_\theta \log p_\theta(t \mid \bm{x})
          \,\nabla_\theta \log p_\theta(t \mid \bm{x})^\top\big]
  \;\succeq\; 0
  \label{eq:fisher}
\end{equation}
is the per-probe Fisher information. Averaging over the probe set,
\begin{equation}
  D_{\mathrm{KL}}(\theta, \theta+\delta)
  \;=\; \tfrac12\, \delta^\top \bar{\bm{F}}_\theta\, \delta
      + O(\norm{\delta}_2^3),
  \qquad
  \bar{\bm{F}}_\theta \;=\; \frac1m \sum_{i=1}^m \bm{F}_\theta(\bm{x}^{(i)})
  \;\succeq\; 0.
  \label{eq:dklexpand}
\end{equation}
Because $\norm{\delta}_2 \propto r$ under the random-perturbation protocol of
\S\ref{sec:exp:calib}, the expected divergence grows \emph{quadratically} in
$r$ for small perturbations and is monotone along any fixed perturbation
direction. By Proposition~\ref{prop:score},
$\nabla_\theta \log p_\theta(t \mid \bm{x})$ is the $t$-th row of
$\bm{J}_\theta(\bm{x})$, so
$\bm{F}_\theta(\bm{x}) = \bm{J}_\theta(\bm{x})^\top
\operatorname{diag}\!\big(p_\theta(\cdot \mid \bm{x})\big)
\bm{J}_\theta(\bm{x})$: the Fisher information is a probability-weighted squared
Jacobian. Because $\bar{\bm{F}}_\theta \succeq 0$, its kernel coincides with the
set of directions in which \emph{every} observed per-token score has zero
projection---the same null space $\mathcal{K}_\theta$ that drives the
adversarial-invariance analysis below.
\end{proposition}

\begin{remark}
\label{rem:linearity}
Quadratic growth in $r$ and monotonicity hold in expectation and in the
small-perturbation regime of Proposition~\ref{prop:fisher}. The monotonicity of
a single realized calibration curve---in which one fixed noise direction is
rescaled by $r$---is an empirical statement, as in \S\ref{sec:exp:calib}.
Because the realized $D_{\mathrm{KL}}$ is, to leading order, the positive
quadratic form $\tfrac12\, r^2 \hat\delta^\top \bar{\bm{F}}_\theta \hat\delta$
with $\bar{\bm{F}}_\theta \succeq 0$, it increases with $r$ along any fixed
direction; the mild sub-quadratic flattening at larger $r$ (e.g.\ layer $8$ in
Table~\ref{tab:calib}) reflects the onset of the saturation regime of
Proposition~\ref{prop:saturation}, not a linear trend.
\end{remark}

\begin{proposition}[Global Lipschitz property and saturation]
\label{prop:saturation}
For \emph{arbitrary} $\delta$ (not necessarily small), the log-probability
vector is a Lipschitz function of the logits: since the log-softmax has
Jacobian $I - \mathbf{1}\,\softmax(\bm{z})^\top$ of spectral norm $1$,
\begin{equation}
  \norm{\bm{\phi}_{\theta+\delta}(\bm{x}) - \bm{\phi}_\theta(\bm{x})}_2
  \;\le\; \norm{\bm{z}_{\theta+\delta}(\bm{x}) - \bm{z}_\theta(\bm{x})}_2
  \;\le\; L_{\mathrm{net}}\,\norm{\delta}_2,
  \label{eq:globlip}
\end{equation}
where $L_{\mathrm{net}}$ is the Lipschitz constant of the logit map in the
weights. In the black-box setting $\bm{\phi}$ is uniformly bounded---each
coordinate is a top-$k$ log-probability ($\le 0$) or a fixed floor---so the
per-probe KL divergence is finite and uniformly bounded in $\theta$. Unlike the
small-scale quadratic regime of Proposition~\ref{prop:fisher}, however, the
divergence does not grow without bound: for large $\delta$ the perturbed
distribution collapses toward a maximally degenerate configuration, and the
ceiling is set by the \emph{clean} model,
\begin{equation}
  D_{\mathrm{KL}}(\theta, \theta+\delta)
  \;\xrightarrow{\ \norm{\delta}_2 \to \infty\ }\;
  \frac1m \sum_{i=1}^m \Big[\log |\Vsh| - H\big(p_\theta(\cdot \mid \bm{x}^{(i)})\big)\Big],
  \label{eq:saturation}
\end{equation}
where $H$ denotes Shannon entropy; each summand is
$\KL\big(p_\theta(\cdot \mid \bm{x}^{(i)}) \,\|\, \mathrm{Uniform}(\Vsh)\big) =
\log |\Vsh| - H(p_\theta(\cdot \mid \bm{x}^{(i)}))$. The ceiling is finite and
independent of further $\delta$; a sufficiently large perturbation is therefore
\emph{indistinguishable from a different model} (the same phenomenon as
\S\ref{sec:exp:cross}).
\end{proposition}

Combining Propositions~\ref{prop:fisher} and \ref{prop:saturation}, the
perturbation response passes through three regimes
(Table~\ref{tab:regimes}). The quadratic regime is valid for
$\norm{\delta}_2 \lesssim \rho = \norm{\bm{J}}/\norm{\bm{H}}$, where $\bm{H}$
is the Hessian of $\bm{\phi}$; here
$D_{\mathrm{KL}} \approx \tfrac12\, r^2 \hat\delta^\top \bar{\bm{F}}_\theta
\hat\delta$, so the divergence scales as $r^2$ with a coefficient set by the
Fisher information. The intermediate regime is polynomial, with higher-order
terms becoming significant; and the saturation regime is a bounded, monotone
flattening toward the entropy ceiling of Eq.~\eqref{eq:saturation}.
Consequently, inverting the monotone calibration curve of
\S\ref{sec:exp:calib}---mapping an observed divergence back to a perturbation
scale---is reliable within the first two regimes, where the response is
strictly increasing in $r$.

\begin{table}[t]
\centering
\small
\caption{Regimes of the fingerprint response to weight perturbation.}
\label{tab:regimes}
\begin{tabular}{lll}
\toprule
Regime & Response & Mechanism \\
\midrule
Small ($\norm{\delta}_2 \lesssim \rho$) & $D_{\mathrm{KL}} \approx c\, r^2$ (quadratic) & Fisher information, Prop.~\ref{prop:fisher} \\
Moderate & $D_{\mathrm{KL}} \approx c\, r^2 + c' r^3 + \cdots$ & higher-order terms, Prop.~\ref{prop:fisher} \\
Large & $D_{\mathrm{KL}} \to \KL(p_\theta \| \mathrm{Uniform}(\Vsh))$ (saturation) & entropy ceiling, Prop.~\ref{prop:saturation} \\
\bottomrule
\end{tabular}
\end{table}

\paragraph{Adversarial invariance.}
We close the theory section with the question a skeptical reader may raise:
whether an adversary who knows the entire method---probe construction and
divergence statistic alike---can engineer a targeted weight edit $\delta \in
\R^P$ (for instance, injecting a backdoor into a served model) that evades
detection. Formally, let
$\bm{F} : \theta \mapsto \big(p_\theta(\cdot \mid \bm{x}^{(1)}), \dots,
p_\theta(\cdot \mid \bm{x}^{(m)})\big)$ denote the fingerprint map, and suppose
the adversary can reconstruct a probe set close to $\mathcal{P}$, has query
access to top-$k$ log-probabilities, and seeks a perturbation $\delta$ that (i)
implements a target edit supported on a ``backdoor block'' $S \subseteq \{1,
\dots, P\}$, with $\norm{\delta_S}_2 \ge b > 0$ for some minimal strength $b$,
and (ii) leaves the fingerprint invariant, $D_{\mathrm{KL}}(\theta,
\theta+\delta) \approx 0$. Stack the probe Jacobians into
\begin{equation}
  \bm{G}_\theta =
  \begin{bmatrix} \bm{J}_1 \\ \vdots \\ \bm{J}_m \end{bmatrix}
  \in \R^{m|\Vsh| \times P},
  \label{eq:gstack}
\end{equation}
so that the $m$ probe distributions---each a vector in $\R^{|\Vsh|}$---are
observed only through the linear map $\delta \mapsto \bm{G}_\theta\delta$.
Directions in $\mathcal{K}_\theta := \ker \bm{G}_\theta$ are exactly the
directions of vanishing quadratic response of the fingerprint-mismatch
objective $\delta \mapsto D_{\mathrm{KL}}(\theta, \theta+\delta)$: by
Proposition~\ref{prop:fisher}, $\ker \bar{\bm{F}}_\theta = \ker \bm{G}_\theta$,
so edits in $\mathcal{K}_\theta$ are invisible to second order by
construction.

\begin{proposition}[Stealth is null-space membership]
\label{prop:stealth}
The KL divergence responds to an edit $\delta$ as
$D_{\mathrm{KL}}(\theta, \theta+\delta) = \tfrac12\, \delta^\top
\bar{\bm{F}}_\theta\, \delta + O(\norm{\delta}_2^3)$
(Proposition~\ref{prop:fisher}). A perturbation is invisible to second
order---the quadratic term vanishes, $\delta^\top \bar{\bm{F}}_\theta\, \delta
= 0$---exactly when $\delta \in \mathcal{K}_\theta = \ker \bm{G}_\theta$, since
$\ker \bar{\bm{F}}_\theta = \ker \bm{G}_\theta$. Moreover, any $\delta \notin
\mathcal{K}_\theta$ carries a \emph{detection floor} that survives the best
possible first-order cancellation: because every observed probability in the
black-box construction is bounded below by a floor $\gamma > 0$, the quadratic
term obeys
\begin{equation}
  \delta^\top \bar{\bm{F}}_\theta\, \delta
  \;=\; \frac1m \sum_{i=1}^m \sum_{t \in \Vsh}
      p_\theta(t \mid \bm{x}^{(i)})\, \big(\bm{J}_{i,t}\,\delta\big)^2
  \;\ge\; \frac{\gamma}{m}\, \norm{\bm{G}_\theta\,\delta}_2^2,
  \label{eq:floor}
\end{equation}
where $\bm{J}_{i,t}$ is the $t$-th row of $\bm{J}_i$; hence
$D_{\mathrm{KL}}(\theta, \theta+\delta) \ge \frac{\gamma}{2m}\,
\norm{\bm{G}_\theta\,\delta}_2^2 - O(\norm{\delta}_2^3)$. The irreducible term
is quadratic in the \emph{visible} part of $\delta$ and invariant under
$\delta \mapsto -\delta$, so no choice of the first-order response removes it.
\end{proposition}

\begin{proposition}[The cost of hiding a visible edit]
\label{prop:comp}
Fix a target edit $\bm{v}_B$ supported on $S$. If $\bm{v}_B \notin
\mathcal{K}_\theta$, any first-order-stealthy perturbation of the form
$\delta = \bm{v}_B + \bm{c}$ must satisfy $\bm{G}_\theta(\bm{v}_B + \bm{c}) =
0$; the minimal-norm compensation is $\bm{c}^\star = -\bm{G}_\theta^{+}
\bm{G}_\theta\,\bm{v}_B = -\Pi_{\mathrm{row}}\,\bm{v}_B$, where
$\Pi_{\mathrm{row}}$ is the orthogonal projection onto the row space of
$\bm{G}_\theta$. Since $\bm{G}_\theta$ annihilates the kernel component of
$\bm{v}_B$, we have $\bm{G}_\theta\bm{v}_B =
\bm{G}_\theta\Pi_{\mathrm{row}}\bm{v}_B$ and hence
\begin{equation}
  \norm{\bm{c}^\star}_2
  \;=\; \norm{\Pi_{\mathrm{row}}\bm{v}_B}_2
  \;\ge\; \frac{\norm{\bm{G}_{\theta,S}\,\bm{v}_B}_2}
                 {\sigma_{\max}\!\big(\bm{G}_\theta\big)},
  \label{eq:compcost}
\end{equation}
where $\bm{G}_{\theta,S}$ restricts $\bm{G}_\theta$ to the columns in $S$ (so
$\bm{G}_\theta\bm{v}_B = \bm{G}_{\theta,S}\bm{v}_B$, since $\bm{v}_B$ is
supported on $S$). Crucially, this bound makes \emph{no rank assumption}: it is
strictly positive exactly when $\bm{G}_{\theta,S}\bm{v}_B \ne 0$, that is,
whenever the target edit has a nonzero \emph{visible} component within its block
($\bm{v}_B \notin \ker \bm{G}_{\theta,S}$), a condition that holds for
Lebesgue-almost-every $\bm{v}_B$ as soon as
$\operatorname{rank}(\bm{G}_{\theta,S}) \ge 1$. In the fully observable case
$\operatorname{rank}(\bm{G}_{\theta,S}) = |S|$ and
$\sigma_{\min}(\bm{G}_{\theta,S}) > 0$, so
$\norm{\bm{G}_{\theta,S}\bm{v}_B}_2 \ge
\sigma_{\min}(\bm{G}_{\theta,S})\,\norm{\bm{v}_B}_2$ and
Eq.~\eqref{eq:compcost} strengthens to the uniform positive fraction of
$\norm{\bm{v}_B}_2$; the rank-free bound is thus a strict generalization of the
rank-based one. In particular, a \emph{support-confined} edit
($\delta = \delta_S$) is stealthy only if $\delta_S \in \ker
\bm{G}_{\theta,S}$, a subspace of dimension $|S| -
\operatorname{rank}(\bm{G}_{\theta,S})$: full rank collapses it to $\{0\}$ (no
nonzero support-confined stealthy edit), while in general it remains a proper,
generically thin subspace, and any edit outside it forces the adversary to
perturb weights outside $S$ with a compensation norm bounded below by its
visible footprint scaled by $1/\sigma_{\max}(\bm{G}_\theta)$---thereby inflating
both the total edit norm and, through Eq.~\eqref{eq:floor}, its second-order
footprint.
\end{proposition}

\begin{remark}[Why evasion is infeasible, not merely hard]
\label{rem:adv}
Propositions~\ref{prop:stealth} and \ref{prop:comp} make precise the sense in
which ``keeping the distributional-divergence fingerprint invariant across
multiple probes is extremely costly.'' The kernel $\mathcal{K}_\theta$ has codimension
$\operatorname{rank}\bm{G}_\theta \le m|\Vsh|$, and since $P$ exceeds
$m|\Vsh|$ by orders of magnitude for the models considered, it is nonempty: the
obstacle is not the absence of invariant directions but their thinness, their
discoverability, and the cost of using them. (i) \emph{Thinness.} The set of
target edits satisfying $\bm{v}_B \in \mathcal{K}_\theta$ is a proper linear
subspace, hence of measure zero in $\R^P$; a backdoor not engineered
specifically for the kernel is visible with certainty. (ii)
\emph{Discoverability.} Engineering it requires the Jacobian. White-box, one
must solve a system of $m|\Vsh|$ equations in $P \sim 10^9$ variables. Because
the loss landscape of an LLM is pathologically ill-conditioned---at the level
of the observable, the singular values of $\bm{G}_\theta$ span many decades; by
Proposition~\ref{prop:score} the Fisher information is $\sum_i \bm{J}_i^\top
\operatorname{diag}\!\big(p_\theta(\cdot\mid \bm{x}^{(i)})\big)\bm{J}_i$, so
the curvature along a weight direction is a probability-weighted squared
observable footprint---any finite-precision solve leaves a residue
$\bm{G}_\theta\hat\delta \ne 0$ that Eq.~\eqref{eq:floor} amplifies into an
irreducible quadratic shift. Black-box, which is the setting this paper
targets, $\bm{G}_\theta$ is not even available: estimating it by forward
differences costs $\Omega(P)$ noisy oracle queries, and any ``null'' direction
$\hat\delta$ of the estimated Jacobian $\hat{\bm{G}}$ satisfies
$\norm{\bm{G}_\theta\hat\delta}_2 \le \norm{\hat{\bm{G}} - \bm{G}_\theta}
\,\norm{\hat\delta}_2$, again leaving a floor for the detector. (iii)
\emph{Cost of use.} Even when an invariant direction exists,
Eq.~\eqref{eq:compcost} shows that hiding a behavior-relevant edit requires a
\emph{global} compensation whose norm is at least the edit's visible footprint
scaled by $1/\sigma_{\max}(\bm{G}_\theta)$, positive whenever the block is not
entirely invisible to the fingerprint (that is, whenever
$\operatorname{rank}(\bm{G}_{\theta,S}) \ge 1$), and the compensation must be
exact: any shortfall is read by the second-order floor.
\end{remark}

\begin{remark}[Quantitative size of the stealth region]
\label{rem:advscale}
The floor of Eq.~\eqref{eq:floor} is quadratic in the visible residue, so the
effective radius of the stealth region shrinks as the fingerprint sensitivity
grows. Under the random-perturbation protocol of \S\ref{sec:exp:calib}
($\delta \sim \mathcal{N}(0, (r\sigma)^2 I)$),
Proposition~\ref{prop:fisher} gives the expected KL shift as
$\tfrac12 (r\sigma)^2 \operatorname{tr}\bar{\bm{F}}_\theta$, and requiring
$D_{\mathrm{KL}} \le \tau$ pins the perturbation scale to
$r\sigma \lesssim \sqrt{2\tau \big/ \operatorname{tr}\bar{\bm{F}}_\theta}$.
The measured sensitivities are large---$D_{\mathrm{KL}}$ reaches
$0.010$--$0.028$ already at $r = 0.02$ and grows as $r^2$
(Table~\ref{tab:calib})---so this radius is small compared with any edit that
materially changes model behavior. The second-order incompressibility of
Proposition~\ref{prop:stealth} is what converts a fingerprint that is merely
\emph{sensitive} into one that is \emph{costly to evade}.
\end{remark}

\subsection{Change Detection}
\label{sec:detect}

Given a reference collection of per-probe next-token distributions from
$M_{\theta_0}$ and a re-probe collected from a candidate model, we combine the
distributional-divergence statistic with a per-token test into a verdict.

\paragraph{Per-probe KL divergence (primary statistic).}
We estimate the per-probe next-token distributions from the observed
log-probabilities (top-$k$ plus floor), and evaluate the mean per-probe KL
divergence $D_{\mathrm{KL}}$ of Eq.~\eqref{eq:dkl}, which by the analysis of
\S\ref{sec:theory} grows quadratically in any weight perturbation and is
exactly zero in the no-change regime.

\paragraph{Per-token two-sample test.}
For each token $t$ observed with real data in both collections, we form the
two-sample $z$-score
\begin{equation}
  z_t = \frac{\mu_{1,t} - \mu_{0,t}}{\sqrt{s_{0,t}^2/n_0 + s_{1,t}^2/n_1}},
  \label{eq:z}
\end{equation}
and sum their squares into a $\chi^2$ statistic with the number of compared
tokens as degrees of freedom; the survival probability $p = \Pr(\chi^2_\nu >
\sum_t z_t^2)$ measures the aggregate mean shift.

\paragraph{Mean-vector and covariance-spectrum diagnostics.}
On the common-token subset, we also report how the distributions moved, using
the mean vector $\bm{\mu}_j = \tfrac1m \sum_{i=1}^m \bar{\bm{\phi}}^{(i)}_j$ of
the per-probe log-probability vectors (collection $j \in \{0,1\}$),
\begin{equation}
  D_{\mathrm{mean}} = \frac{\norm{\bm{\mu}_1 - \bm{\mu}_0}}{\norm{\bm{\mu}_0}} ,
  \label{eq:dmean}
\end{equation}
and their covariance spectra $\bm{\lambda}_0, \bm{\lambda}_1$,
\begin{equation}
  D_{\mathrm{cov}} = \frac{\norm{\bm{\lambda}_1 - \bm{\lambda}_0}}{\norm{\bm{\lambda}_0}} .
  \label{eq:dcov}
\end{equation}
These are diagnostic only: $D_{\mathrm{mean}}$ is sensitive but weak, and
$D_{\mathrm{cov}}$ can saturate or even fall at large perturbations, so neither
serves as the primary statistic (\S\ref{sec:exp:calib}).

A model is declared \textsc{changed} if the per-token $\chi^2$ test is
significant at level $\alpha$; otherwise \textsc{unchanged}. The mean per-probe
KL divergence $D_{\mathrm{KL}}$ is reported alongside the verdict as the
interpretable magnitude of the change. In the zero-noise regime the re-probe
reproduces the reference distributions exactly, so $D_{\mathrm{KL}} = 0$ and
the detector admits no false positives (\S\ref{sec:exp}).

\section{Experiments}
\label{sec:exp}

\paragraph{Setup.}
Local calibration experiments use three small open-weight models:
\texttt{Qwen\slash Qwen3.5-0.8B} ($|\V|=248{,}077$ subword tokens),
\texttt{unsloth\slash Llama-3.2-1B} ($|\V|=128{,}256$), and
\texttt{unsloth\slash gemma-3-1b-pt} ($|\V|=262{,}145$), at intermediate
layer $8$ (and layer $16$ for Qwen) with mean pooling; the black-box
experiments use the DeepSeek API (\texttt{deepseek-v4-flash}, $|\V|=129{,}280$)
which exposes
top-$20$ log-probabilities with reasoning disabled. The shared subword
vocabulary with DeepSeek contains $108{,}773$ tokens for Qwen ($43.8\%$ of its
vocabulary), $70{,}663$ for Llama, and only $21{,}237$ for Gemma ($8.1\%$ of
its much larger vocabulary), so the \emph{surface} overlap varies
substantially across families. As \S\ref{sec:exp:cross} shows, however, even a
large surface overlap does not translate into prediction overlap.

\subsection{Probe Quality and the Limits of Alignment}
\label{sec:exp:syn}

\begin{figure}[t]
\centering
\includegraphics[width=0.48\linewidth]{figures/fig_discriminative.pdf}
\caption{Discriminative power of hidden-state cosine similarity, by layer.
Representation-equivalent (eq) pairs and unrelated (rnd) pairs are compared;
the annotated gap (mean cosine difference) peaks at layer $16$.}
\label{fig:disc}
\end{figure}

Our first result is negative but informative: intermediate-layer alignment
alone is a \emph{necessary but not sufficient} condition for being
representation-equivalent. Sequences that align closely with $\bm{h}^\star$ are
not always fluent or semantically coherent---the latent-cluster objective
admits degenerate solutions; with an unconstrained candidate set the search collapses to
the reference itself unchanged. Quantitatively, hidden-state cosine similarity
carries only a weak discriminative signal: across layers, the cosine gap between
representation-equivalent (eq) and unrelated (rnd) sentence pairs is
$0.18$--$0.32$, with equivalent pairs at cosine $\approx0.85$--$0.95$ and
unrelated pairs at $\approx0.60$--$0.73$ (Figure~\ref{fig:disc}). Restricting
the candidate set to embedding-space neighbors within $\Vsh$
(\S\ref{sec:syn}) resolves the fluency problem practically, yielding
representation-equivalent pairs such as \emph{``the lazy dog'' $\to$ ``the lazy
cat''} while preserving alignment (mean cosine $\approx0.98$). This pair is
\emph{not} a synonym pair in the linguistic sense---the two sentences differ in
meaning to a human reader---yet the model's representation places them in a
tight cluster, which is exactly the notion of equivalence this paper relies on.

\subsection{Calibration under Gaussian Weight Noise}
\label{sec:exp:calib}

\begin{table}[t]
\centering
\small
\caption{Weight-noise calibration under the shared-vocabulary restriction:
mean per-probe next-token KL divergence $D_{\mathrm{KL}}$ as a function of
injected noise scale $r$ (Gaussian $N(0,(r\sigma)^2)$ per tensor), for the four
open-weight configurations (log-probabilities truncated to a floor of $-30$).}
\label{tab:calib}
\begin{tabular}{llc}
\toprule
Model & $r$ & $D_{\mathrm{KL}}$ \\
\midrule
\multirow{8}{*}{Qwen-0.8B (L8)} & 0.00 & 0.0000 \\
 & 0.02 & 0.0164 \\
 & 0.05 & 0.0992 \\
 & 0.10 & 0.5246 \\
 & 0.20 & 3.2754 \\
 & 0.40 & 9.2107 \\
 & 0.60 & 12.9029 \\
 & 1.00 & 14.7039 \\
\midrule
\multirow{8}{*}{Qwen-0.8B (L16)} & 0.00 & 0.0000 \\
 & 0.02 & 0.0191 \\
 & 0.05 & 0.1216 \\
 & 0.10 & 0.6119 \\
 & 0.20 & 3.4085 \\
 & 0.40 & 8.8989 \\
 & 0.60 & 12.5548 \\
 & 1.00 & 14.8757 \\
\midrule
\multirow{8}{*}{Llama-3.2-1B} & 0.00 & 0.0000 \\
 & 0.02 & 0.0102 \\
 & 0.05 & 0.0562 \\
 & 0.10 & 0.2181 \\
 & 0.20 & 1.0944 \\
 & 0.40 & 5.7986 \\
 & 0.60 & 9.6807 \\
 & 1.00 & 11.3151 \\
\midrule
\multirow{8}{*}{Gemma-3-1b} & 0.00 & 0.0000 \\
 & 0.02 & 0.0275 \\
 & 0.05 & 0.2301 \\
 & 0.10 & 0.9680 \\
 & 0.20 & 4.7647 \\
 & 0.40 & 13.1882 \\
 & 0.60 & 23.0286 \\
 & 1.00 & 26.6323 \\
\bottomrule
\end{tabular}
\end{table}

\begin{figure}[t]
\centering
\includegraphics[width=0.95\linewidth]{figures/fig_calibration.pdf}
\caption{Weight-noise calibration for the open-weight configurations. The mean
per-probe next-token KL divergence $D_{\mathrm{KL}}$ grows monotonically in the
noise scale $r$---quadratically at small $r$ (matching the Fisher-information
prediction) and saturating at large $r$; on a log--log axis the small-scale
regime appears as a slope-$2$ line. Dashed lines are quadratic fits in the
small-scale regime.}
\label{fig:calib}
\end{figure}

To validate that the distributional-divergence fingerprint tracks the weights,
we inject zero-mean Gaussian noise $N(0, (r\sigma)^2)$ into the weights of each
model, where $\sigma$ is the per-tensor parameter scale, and measure how the
mean per-probe next-token KL divergence from the clean model behaves as a
function of $r$. Table~\ref{tab:calib} and Figure~\ref{fig:calib} report the
results for the four open-weight configurations. In every case $D_{\mathrm{KL}}$
grows monotonically in $r$, and within the small-scale regime it grows
approximately \emph{quadratically}: increasing $r$ from $0.02$ to $0.1$ (a
factor of $5$) raises $D_{\mathrm{KL}}$ by roughly a factor of $30$ (e.g.\ Qwen,
layer $8$: $0.016 \to 0.525$), matching the Fisher-information prediction
$D_{\mathrm{KL}} \approx \tfrac12 \delta^\top \bar{\bm{F}}_\theta \delta \propto
r^2$ of \S\ref{sec:theory}. At larger $r$ the growth slows toward the entropy
ceiling of Proposition~\ref{prop:saturation}.

We chose $D_{\mathrm{KL}}$ after ablating the alternative statistics on the
same perturbation protocol. The mean-vector divergence $D_{\mathrm{mean}}$ is
monotonic but weak---it reaches only $0.545$ (Qwen, layer $8$) and $0.392$
(Llama) at $r = 1.0$, against $14.70$ and $11.32$ for $D_{\mathrm{KL}}$. The
covariance-spectrum divergence $D_{\mathrm{cov}}$ is more sensitive at small
$r$ but \emph{non-monotonic}: its spectrum diverges at moderate $r$ and then
collapses as the perturbed distribution concentrates (Qwen: $0.046 \to \cdots
\to 0.659 \to 0.592$; Llama: $0.012 \to \cdots \to 0.841 \to 0.661$). The KL
divergence alone combines monotonicity with sensitivity, which is why it serves
as the primary statistic in \S\ref{sec:detect}; $D_{\mathrm{mean}}$ and
$D_{\mathrm{cov}}$ are retained only as diagnostics. This monotone, quadratic
small-scale calibration is the core evidence that the fingerprint is a
faithful, sensitive proxy for weight changes rather than a measure of surface
agreement: inverting the monotone calibration curve maps an observed divergence
back to the implied relative perturbation scale $r$.

\subsection{Zero-Noise Change Detection over the API}
\label{sec:exp:api}

Using the DeepSeek API, we collect a reference distributional fingerprint over
a probe set of $15$ representation-equivalent contexts (two samples each), then
re-probe the same contexts.
Because the API is deterministic under fixed parameters, the re-probe
reproduces the reference statistics exactly: the per-token $\chi^2$ test is
null ($\chi^2 = 0.00$, $df = 28$, $p = 1$), and the mean per-probe KL divergence
is $D_{\mathrm{KL}} = 0.000$, yielding an \textsc{unchanged} verdict with no
false positives. The recorded log-probability field is highly concentrated: it
has trace $\operatorname{tr}(\bm{\Sigma}) = 516.8$ over $60$ observed tokens,
with the top three eigenvalues accounting for $73\%$ of the total variance---a
low-rank structure (Figure~\ref{fig:spectrum}) that reflects the fact that the
probes all collapse onto a single representation-equivalence cluster. This
spectral concentration is a characterization of the probe set rather than the
primary statistic; the change verdict rests on the divergence. The fingerprint
therefore acts as a zero-noise change detector suitable for auditing a served
model over time.

\begin{figure}[t]
\centering
\includegraphics[width=0.48\linewidth]{figures/fig_spectrum.pdf}
\caption{Spectrum of the per-probe log-probability covariance of the DeepSeek
fingerprint (log scale). The rapid decay of eigenvalues, with the top three
accounting for $73\%$ of the trace, reflects the concentration of
representation-equivalent probes onto a single cluster.}
\label{fig:spectrum}
\end{figure}

\subsection{Cross-Model Comparison}
\label{sec:exp:cross}

Comparing a Qwen checkpoint against the DeepSeek model directly is frustrated
by a subtle form of mismatch. Although the two families share $108{,}773$
subword tokens, their \emph{predictions} barely overlap: on a shared context,
the top-$20$ next-tokens predicted by the two models have an average Jaccard
overlap of only $0.012$ (on average $0.7$ shared tokens, out of roughly $39$
and $17$ predicted tokens respectively). The representation-equivalent probes therefore
induce log-probability vectors whose supports are almost disjoint across the
two models, so the vectors are not directly comparable. This is a structural
limitation of distribution-level fingerprinting across distinct model
families, and is distinct from the \emph{within}-family change detection in
which the model is fixed and only its weights are perturbed.

\section{Discussion and Limitations}
\label{sec:disc}

\paragraph{What the distributional divergence buys over the mean and covariance.}
The mean log-probability vector is a first-order summary that is dominated by
the model's coarse language statistics and is relatively insensitive to the
precise weights. The covariance of the per-probe log-probability vectors is
more sensitive, but its response is dominated by the spectrum and can saturate
or collapse at large perturbations. The per-probe distributional divergence
(KL) combines the best of both: it is far more sensitive than the mean---its
quadratic response is governed by the Fisher information, which accumulates
probability-weighted squared Jacobians rather than averaging them away---and it
remains monotone where the covariance saturates, because the KL response is a
positive quadratic form for small perturbations and a finite, entropy-governed
ceiling for large ones.

\paragraph{Limitations.}
(i) The method is fundamentally dependent on the service provider exposing
token-level log-probabilities: if an API only returns a final answer---or
quantizes or withholds its top-$k$ log-probabilities---the fingerprint cannot
be extracted. (ii) Latent alignment is necessary but not
sufficient for representation-equivalence, so probe quality depends on the
candidate-restriction heuristic. (iii) The fingerprint
is tokenizer-relative, so cross-family comparison remains open. (iv) The
top-$k$ truncation of the API injects a log-probability floor (here $-30$)
below which the next-token distribution is only bounded, not measured; this
floors the estimated $D_{\mathrm{KL}}$ and should be accounted for when setting
detection thresholds.
(v) Our calibration is demonstrated at a single model scale ($0.8$B); whether
the quadratic (Fisher-information) regime holds at much larger scales is
unverified. (vi) An adversary
who controls the served model and knows the method is constrained but not
precluded by the analysis of \S\ref{sec:theory}: second-order stealth is
exactly null-space membership (Prop.~\ref{prop:stealth}), so evasion is
possible in principle---any distribution-level detector inherits the
output-equivalence ambiguity of \S\ref{sec:intro}---but it requires engineering
a weight edit into a Jacobian-defined, measure-zero subspace whose discovery is
query- and precision-bound, at a compensation cost bounded below by
Eq.~\eqref{eq:compcost}. Randomized or secret probe sets and robust-statistic
variants remain prudent hardening. (vii) The monotone, invertible quadratic (Fisher-information)
regime of \S\ref{sec:exp:calib} has a finite radius: beyond the
small-perturbation scale the fingerprint saturates toward the entropy ceiling
(\S\ref{sec:theory}, Prop.~\ref{prop:saturation}), and a sufficiently large
perturbation is indistinguishable from a different model---a blind spot
analogous to the cross-family barrier of \S\ref{sec:exp:cross}.

\section{Conclusion}
\label{sec:concl}

We have shown that the per-probe distributional divergence (KL) of next-token
distributions, induced by actively synthesized representation-equivalent
probes, is a practical and sensitive fingerprint of a language model's weights.
The approach is black-box-compatible (requiring only top-$k$
log-probabilities), yields a monotonic calibration curve under injected
Gaussian weight noise, and admits zero-noise change detection over a production
API. The central finding is conceptual: to detect weight equivalence, one
should compare the \emph{second-order} response of the observable next-token
distribution rather than its expectation---a response governed by the Fisher
information---and representation-equivalent probes---sentences the model itself
cannot tell apart---provide the controlled setting in which that response is
identifiable.

\paragraph{Acknowledgments.} We thank the anonymous reviewers for their
feedback.

\bibliographystyle{plainnat}
\bibliography{refs}

\end{document}
